
    
\subsection{Few types in balls}\label{sec:few-types-in-balls}
In this subsection we prove \Cref{lem:ball-types}.
We use the following Ramsey-type result for set systems due to Ding, Oporowski, Oxley, and Vertigan~\cite{DBLP:journals/jct/DingOOV96}.

% lax begin Lax185574.AtomicMatching
\begin{definition}
    For a graph $G$ and tuples $\bar a, \bar b \in V(G)^t$, we say $\bar a = a_1 \ldots a_t$ and $\bar b = b_1 \ldots b_t$ are \emph{atomically matched} 
    if there is a relation ${\sim} \in \{=,\neq,\geq\}$ such that $G \models E(a_i, b_j)$ if and only if $i \sim j$ for all $i,j \in [t]$.
    In particular, the bipartite graph semi-induced between the tuples $\bar a$ and $\bar b$ is a matching, a co-matching, or a half-graph.
\end{definition}
% lax end

% lax begin Lax185574.MatchedTuples
% lax begin Lax185574Proofs.MatchedTuples.exists_atomicallyMatched
\begin{theorem}[{\cite[Corollary~2.4]{DBLP:journals/jct/DingOOV96}, see also~\cite[Theorem~2]{gravier2004}}]\label{thm:matching}
    There exists a computable function $Q:\N\rightarrow\N$ such that for every $\ell\in\N$ and for every bipartite graph $G = (L,R,E)$ with
    $|\ATypes(L/R)| > Q(\ell)$, there are two tuples $\bar a \in L^\ell$ and $\bar b \in R^\ell$ that are atomically matched in $G$.
\end{theorem}
% lax end
% lax end

We use the above result to construct the following obstructions to monadic dependence.

% lax begin Lax185574.PhiPsiPatterns
\begin{restatable}{definition}{defphipsi}\label{def:phi-psi-patterns}
    Fix first-order formulas $\phi(\bar x,z)$ and $\psi(\bar y,z)$. A \emph{$(\phi,\psi)$-pattern} of order $s \times t$ in a graph~$G$ consists of tuples $\bar a_1, \ldots, \bar a_s \in V(G)^{|\bar x|}$, $\bar b_1, \ldots, \bar b_t \in V(G)^{|\bar y|}$, and vertices ${(c_{i,j}: i\in [s],j\in [t])}$ such that the following holds
    for all $i,i'\in [s]$ and $j,j'\in[t]$:
    \begin{enumerate}
        \item $G\models \phi(\bar a_i,c_{i,j})$ 
        and $G\not\models \phi(\bar a_{i'},c_{i,j})$ if $i'>i$,
        \item $G\models \psi(\bar b_j,c_{i,j})$ 
        and $G\not\models \psi(\bar b_{j'},c_{i,j})$ if $j'>j$.
    \end{enumerate}
\end{restatable}
% lax end
\noindent We call $\setof{c_{i,j}}{i\in[s],j\in[t]}$ the \emph{grid-points} of the pattern, and $s$ its \emph{height}. See \Cref{fig:phi-psi-pattern}.
\begin{figure}[htbp]
    \definecolor{patternphi}{RGB}{0,90,200}
    \definecolor{patternpsi}{RGB}{200,40,40}
    \centering
    \begin{tikzpicture}[
        x=1.05cm,y=0.90cm,font=\small,
        grid point/.style={circle,fill=black,inner sep=1.5pt},
        % Tuples are drawn as o=o: two grid-point-sized dots (radius
        % 1.5pt*sqrt(2)) joined by a bar.
        a tuple/.style={rectangle,inner sep=0pt,
            minimum width=10pt,minimum height=5pt,
            path picture={
                \fill ([xshift=2.12pt]path picture bounding box.west) circle (2.12pt);
                \fill ([xshift=-2.12pt]path picture bounding box.east) circle (2.12pt);
                \fill ([xshift=2.12pt,yshift=-2.12pt]path picture bounding box.west)
                    rectangle ([xshift=-2.12pt,yshift=2.12pt]path picture bounding box.east);
            }},
        b tuple/.style={rectangle,inner sep=0pt,
            minimum width=5pt,minimum height=10pt,
            path picture={
                \fill ([yshift=-2.12pt]path picture bounding box.north) circle (2.12pt);
                \fill ([yshift=2.12pt]path picture bounding box.south) circle (2.12pt);
                \fill ([xshift=-2.12pt,yshift=-2.12pt]path picture bounding box.north)
                    rectangle ([xshift=2.12pt,yshift=2.12pt]path picture bounding box.south);
            }},
        phi incidence/.style={draw=patternphi,line width=0.5pt},
        psi incidence/.style={draw=patternpsi,line width=0.5pt}
    ]
        \foreach \i in {1,...,5} {
            \node[a tuple] (a-\i) at (-0.2,-\i) {};
        }
        \foreach \i/\aindex in {1/1,5/s} {
            \node[anchor=east,text=patternphi,inner sep=1pt] at ([xshift=-1pt]a-\i.west) {$\bar a_{\aindex}$};
        }
        \node[anchor=east,text=patternphi,inner sep=1pt] at ([xshift=-1pt]a-3.west) {$\bar a_i$};
        \foreach \j in {1,...,5} {
            \node[b tuple] (b-\j) at (\j,0.15) {};
        }
        \foreach \j/\bindex in {1/1,5/t} {
            \node[text=patternpsi] at (\j,0.65) {$\bar b_{\bindex}$};
        }
        \node[text=patternpsi] at (3,0.65) {$\bar b_j$};
        \foreach \i in {1,...,5} {
            \foreach \j in {1,...,5} {
                \coordinate (c-\i-\j) at (\j,-\i);
            }
        }
        % Curved fans separate the individual required formula incidences.
        \foreach \i in {1,...,5} {
            \foreach \j in {1,...,5} {
                \pgfmathsetmacro{\incidencewidth}{0.5+0.7*(\i==3)*(\j==3)}
                \draw[phi incidence,line width=\incidencewidth pt] (a-\i.east)
                    .. controls (0.45,-\i+0.05*\j)
                    and (\j-0.25,-\i+0.10+0.09*\j) .. (c-\i-\j);
                \draw[psi incidence,line width=\incidencewidth pt] (b-\j.south)
                    .. controls (\j-0.12-0.10*\i,-0.35)
                    and (\j-0.12-0.10*\i,-\i+0.35) .. (c-\i-\j);
            }
        }
        % Four illustrative phi-incidences with earlier row tuples.
        \foreach \earlier/\i/\j in {2/3/1,3/5/2,4/5/4} {
            \pgfmathsetmacro{\incidencewidth}{0.5+0.7*(\i==3)*(\j==3)}
            \draw[phi incidence,dashed,line width=\incidencewidth pt] (a-\earlier.east)
                .. controls (0.55,-\earlier)
                and (\j-0.40,-\i+0.40) .. (c-\i-\j);
        }
        \draw[phi incidence,dashed,line width=1.2pt] (a-2.east)
            .. controls (0.55,-1.75) and (2.60,-2.60) .. (c-3-3);
        % Five illustrative psi-incidences with earlier column tuples.
        \foreach \earlier/\i/\j in {1/1/2,1/2/3,2/1/3,4/2/5} {
            \pgfmathsetmacro{\incidencewidth}{0.5+0.7*(\i==3)*(\j==3)}
            \draw[psi incidence,dashed,line width=\incidencewidth pt] (b-\earlier.south)
                .. controls (\earlier-0.45,-0.40)
                and (\j-0.35,-\i+0.35) .. (c-\i-\j);
        }
        \draw[psi incidence,dashed] (b-1.south)
            .. controls (0.35,-0.65) and (1.65,-2.90) .. (c-3-2);
        \foreach \i in {1,...,5} {
            \foreach \j in {1,...,5} {
                \node[grid point] at (c-\i-\j) {};
            }
        }
        \foreach \i/\j/\rowindex/\colindex in {1/1/1/1,1/5/1/t,5/1/s/1,5/5/s/t} {
            \node[anchor=north west,fill=white,inner sep=1pt,font=\scriptsize]
                at (\j+0.10,-\i-0.17) {$c_{\rowindex,\colindex}$};
        }
        \node[anchor=north west,fill=white,inner sep=1pt,font=\scriptsize]
            at (3.10,-3.17) {$c_{i,j}$};
        \node[grid point,inner sep=2.8pt] at (c-3-3) {};

        % Truth profiles at the marked grid-point; ? means unrestricted.
        \node[anchor=west] at (6.0,-0.3) {At a grid-point $c_{i,j}$:};
        \node[anchor=west,text=patternphi] at (6.0,-1.1)
            {$\phi(\bar x,c_{i,j})$};
        \node[anchor=west,text=patternpsi] at (6.0,-3.3)
            {$\psi(\bar y,c_{i,j})$};
        \foreach \k/\profilex/\p/\q in {1/7/?/?,2/8.35/?/?,3/9/1/1,4/9.65/0/0,5/11/0/0} {
            \node[draw,minimum size=5mm,text=patternphi] (phi-box-\k)
                at (\profilex,-2.1) {$\p$};
            \node[draw,minimum size=5mm,text=patternpsi] (psi-box-\k)
                at (\profilex,-4.3) {$\q$};
        }
        \foreach \k/\aindex/\bindex in {1/1/1,2/i-1/j-1,3/i/j,4/i+1/j+1,5/s/t} {
            \node[font=\scriptsize] (phi-label-\k)
                at ($(phi-box-\k.center)+(0,0.45)$) {$\bar a_{\aindex}$};
            \node[font=\scriptsize,anchor=south] (psi-label-\k)
                at ($(phi-label-\k.south)+(0,-2.2)$) {$\bar b_{\bindex}$};
        }
        \foreach \leftindex/\rightindex in {1/2,4/5} {
            \node[font=\scriptsize]
                at ($(phi-box-\leftindex.center)!0.5!(phi-box-\rightindex.center)$) {$\cdots$};
            \node[font=\scriptsize]
                at ($(psi-box-\leftindex.center)!0.5!(psi-box-\rightindex.center)$) {$\cdots$};
        }
    \end{tikzpicture}
    \caption{A $(\phi,\psi)$-pattern of order $s\times t$ for $s=t=5$.
        % Horizontal ellipses represent the tuples $\bar a_i$, and vertical ellipses represent the tuples $\bar b_j$.
        Edges indicate $\phi(\bar a_{i'},c_{i,j})$ or $\psi(\bar b_{j'},c_{i,j})$.
        Each grid-point $c_{i,j}$ is related via $\phi$ to the tuple $\bar a_i$ in its row and possibly (indicated by dashed edges) to some earlier row tuples, and analogously for $\psi$ and the column tuples.
        % For each $c_{i,j}$, the truth values at later indices $i'>i$ and $j'>j$
        % must be $0$, while those at earlier indices are unrestricted, denoted by~$?$.
        The profiles on the right illustrate these conditions at a grid-point $c_{i,j}$.}
    \label{fig:phi-psi-pattern}
\end{figure}



\medskip
The following lemma is a corollary of~\cite[Lemma~C.1]{ordered-tww}. We include a self-contained proof in \Cref{app:patterns} for completeness.

% lax begin Lax185574.ForbiddenPhiPsiPatterns
\begin{restatable}{lemma}{lemphipsi}\label{lem:phi-psi-patterns}
    For every monadically dependent graph class $\CC$ and formulas $\phi(\bar x,z),\psi(\bar y,z)$, there is some $t\in\N$ such that no graph $G\in \CC$ contains a $(\phi,\psi)$-pattern of order $t \times t$.
\end{restatable}
% lax end



% lax begin Lax185574.PhiPatterns
\begin{definition}
    For a first-order formula $\phi(\bar x,z)$, a \emph{$\phi$-pattern} is a $(\phi,\psi)$-pattern
    with $\psi(y,z)=E(y,z)$ or $\psi(y,z)=\neg E(y,z)$.
    We call those a
    \emph{$(\phi,E)$-pattern} 
    and a \emph{$(\phi,\neg E)$-pattern}, respectively.
    Such a pattern consists of tuples $\bar a_1,\ldots,\bar a_s\in V(G)^{|\bar x|}$,
    vertices $b_1,\ldots,b_t\in V(G)$ and $(c_{i,j}:i\in[s],j\in [t])$.
    We call $\bar b=(b_1,\ldots,b_t)$
    the \emph{base} of the pattern. 
\end{definition}
% lax end

% lax begin Lax185574.ForbiddenPhiPatterns
% lax begin Lax185574Proofs.ForbiddenPhiPatterns.forbidden_phi_patterns
\begin{corollary}\label{lem:forbidden-patterns}
        For every monadically dependent graph class $\CC$ and formula $\phi(\bar x,z)$, there is some $t\in \N$ such that no graph $G\in \CC$ contains a $\phi$-pattern of order $t \times t$.
\end{corollary}
% lax end
% lax end


We now prove \Cref{lem:ball-types}, which we repeat for convenience.

% lax begin Lax185574.BallTypes
\lemBallTypes*
% lax end
% lax begin Lax185574Proofs.BallTypes.ball_types
\begin{proof}
    For each symmetric relation $R\subseteq\Pi^s \times \Pi^s$, let $\phi_R(\bar z x,y)$ be a first-order formula expressing that $x$ and $y$ have distance at most $r$ in the flip $G\oplus(\bar z,R)$, where $\bar z$ is an $s$-tuple.
Its existence follows from \Cref{obs:sim-def-flip}.

Let $k\ge 1$ be the bound from 
     \Cref{lem:forbidden-patterns}, such that no $G\in\CC$ contains a $\phi_R$-pattern of order $k\times k$, for any $R$.
    Let $Q\from\N\to\N$ be as in \Cref{thm:matching}, and set
    \[
        t\coloneqq Q\big(k \cdot |\Pi^s|^2\big)=Q(k\cdot 9^s).
    \]

    Fix $G\in\CC$ and a finite set $A\subseteq V(G)$.
    For every $R$ as above and $N\in \set{E,\neg E}$, and every $\bar b\in A^k$, choose a maximum-height $(\phi_R,N)$-pattern with base $\bar b$ in $G$. Each of them is an $\ell\times k$ pattern for some $\ell<k$ (possibly $\ell=0$).
    
    Let $A^*$ consist of $A$ and all grid-points of the chosen patterns.
    There are $O_s(1)$ many choices of $R$ and $|A|^k$ many choices of $\bar b$, and each chosen pattern contributes at most $k^2$ vertices, so
    \[
        |A^*|\leq |A|+O_s(1)\cdot|A|^k\cdot k^2
        =O_{\CC,r,s}(|A|^k).
    \]

    Fix a symmetric relation $R\subseteq\Pi^s \times \Pi^s$, tuple $\bar p \in V(G)^s$, and vertex $v \in V(G)$. Let
    $H=G\oplus(\bar p,R)$ and suppose $B\coloneqq \Ball_r^H(v)$ satisfies $B\cap A^*=\emptyset$.
    We show that $B$ realizes at most $t$ atomic types over $A$ in $H$, thus proving the lemma.
    
    Suppose the contrary.
    By \Cref{thm:matching}, there are tuples $\bar u'$ in $B$ and $\bar w'$ in $A$, each of length $k \cdot |\Pi^s|^2$, that are atomically matched in $H$.

    Let $\PP$ be the partition of $V(G)$ into atomic types over $\bar p$; then $H$ is a $\PP$-flip of $G$ and $|\PP|\le |\Pi^s|$.
    By the pigeonhole principle, we can pick corresponding subsequences $\bar u$ of $\bar u'$ and $\bar w$ of $\bar w'$ of length $k$,
    such that $\bar u$ 
    is contained in a single part $X\in \PP$
    and $\bar w$ is contained in a single part $Y\in\PP$.
    Write $\bar u=(u_1,\ldots,u_k)$ and $\bar w=(w_1,\ldots,w_k)$.
    Since $H$ is a $\PP$-flip of $G$, adjacency between $X$ and $Y$ is either preserved or complemented.
    Together with the matching condition in $H$, this gives a choice of $N\in\set{E,\neg E}$ such that
    \[
        \begin{aligned}
            G&\models N(w_j,u_j)&\text{for every }j\in[k],\\
            G&\not\models N(w_{j'},u_j)&\text{for }1\le j<j'\le k.
        \end{aligned}
    \]
    
    Consider the maximum-height $(\phi_R,N)$-pattern chosen for $R$ and $N$ and for  base $\bar b\coloneqq\bar w=(b_1,\ldots,b_k)$.
    Write $\bar a_1,\ldots,\bar a_\ell$ for its witnessing tuples and $c_{i,j}$ for its grid-points, where $i\in[\ell]$ and $j\in[k]$ (possibly $\ell=0$).
    All these grid-points belong to $A^*$.
    Append a new row by setting
    \[
        \bar a_{\ell+1}\coloneqq\bar p v,
        \qquad
        c_{\ell+1,j}\coloneqq u_j\quad\text{for every }j\in[k].
    \]
    We verify that this yields a $(\phi_R,N)$-pattern with the same base $\bar b$ and height $\ell+1$, contradicting the maximality of $\ell$.        
    The column condition holds by the choice of $N$, $b_j=w_j$ and $c_{\ell+1,j}=u_j$:
    \[
        \begin{aligned}
            G&\models N(b_j,c_{\ell+1,j})&\text{for every }j\in[k],\\
            G&\not\models N(b_{j'},c_{\ell+1,j})&\text{for }1\le j<j'\le k.
        \end{aligned}
    \]    
    For every $j\in[k]$, we have $G\models\phi_R(\bar a_{\ell+1},c_{\ell+1,j})$, since $c_{\ell+1,j}\in B=\Ball_r^H(v)$.
    Moreover, for every $i\in[\ell]$, we have $G\not\models\phi_R(\bar a_{\ell+1},c_{i,j})$, since $c_{i,j}\in A^*$ and $B\cap A^*=\emptyset$. This proves the lemma.
\end{proof}
% lax end



